\section{The cubical block, exact quotient, and retained map}\label{sec:block}

Let $I_2$ be the cellular interval with vertices $\{0,1,2\}$ and two positively
oriented edges.  Set
\[
 B=I_2^4.
\]
Write $B_q$ for its oriented $q$-cells and $C^q(B;\R)$ for real cellular
cochains with the coordinate Euclidean inner product.  The coboundaries are
\[
 d_0:C^0(B;\R)\longrightarrow C^1(B;\R),\qquad
 d_1:C^1(B;\R)\longrightarrow C^2(B;\R).
\]
Product counting gives
\begin{equation}\label{eq:cell-counts}
 |B_0|=81,\qquad |B_1|=216,\qquad |B_2|=216.
\end{equation}

Assign $U_e\in\SUtwo$ to each oriented link, with $U_{\bar e}=U_e^{-1}$.
For a plaquette $p=(x;\mu,\nu)$, $\mu<\nu$, put
\begin{equation}\label{eq:plaquette-holonomy}
 U_p=U_{x,\mu}U_{x+\hat\mu,\nu}
 U_{x+\hat\nu,\mu}^{-1}U_{x,\nu}^{-1}.
\end{equation}
The gauge group $\SUtwo^{B_0}$ acts by
$U_{xy}\mapsto g_xU_{xy}g_y^{-1}$.

\subsection{Rooted tree gauge}

Fix the root $(0,0,0,0)$ and the lexicographic spanning tree $T$ selected by
Kruskal's algorithm from the lexicographically ordered positive links.  We
divide by the based gauge group
\[
 \mathcal G_0=\{g\in\SUtwo^{B_0}:g_{(0,0,0,0)}=1\}.
\]
Every orbit has a unique representative satisfying $U_e=1$ for $e\in T$.
There are $80$ tree links and $136$ cotree links.  Successive use of left and
right invariance of Haar measure gives the exact factorization
\begin{equation}\label{eq:haar-factorization}
 \prod_{e\in B_1}\Haar_e
 =\left(\prod_{x\ne(0,0,0,0)}\Haar_{g_x}\right)
   \left(\prod_{e\notin T}\Haar_e\right).
\end{equation}
The proof, including uniqueness, is in Appendix~\ref{app:tree-gauge}.  Thus
the reduced measure has no omitted gauge determinant.  Simultaneous
conjugation by a root group element remains as a global symmetry rather than a
redundancy of the based quotient.

On the principal exponential chart write $U_e=\exp A_e$ for $e\notin T$ and
$A_e=0$ for $e\in T$, with $A_e\in\suTwo\cong\R^3$.  In these coordinates
\begin{equation}\label{eq:haar-log-density}
 \Haar_e=c\,j(A_e)\dd A_e,\qquad
 j(A)=\left(\frac{\sin\norm A}{\norm A}\right)^2,
\end{equation}
where $c$ is constant and the value at the origin is defined by continuity.

\subsection{Retained face-curvature coordinates}

For each pair $\mu<\nu$ and each choice of the other two coordinates in
$\{0,2\}$, sum $d_1A$ over the four fine plaquettes in the corresponding
$2\times2$ boundary face.  This defines a scalar map
\[
 C:C^1(B;\R)\longrightarrow\R^{24}.
\]
It vanishes on $\operatorname{im}d_0$.  Seven relations reduce its rank to
$17$.  The supplement selects the following zero-based row indices:
\begin{equation}\label{eq:selected-rows}
 9,6,10,7,11,0,4,8,1,5,2,17,15,14,16,13,20.
\end{equation}
Let $R:\R^{136}\to\R^{17}$ be these rows restricted to the cotree links.
The three-color retained map is $R\otimes\Id_3$.  It is covariant under the
remaining global conjugation because the logarithm intertwines conjugation
with the adjoint representation inside the principal chart.

\begin{definition}[Retained and fiber spaces]\label{def:retained-fiber}
The scalar retained space is $\mathcal Y_0=\operatorname{ran}R$ and the scalar
fiber is $F_0=\ker R$.  Their $\suTwo$ versions are
\[
 \cY=\mathcal Y_0\otimes\R^3,\qquad
 \cF=F_0\otimes\R^3.
\]
\end{definition}

The retained coordinate is deliberately linear in the gauge-fixed logarithmic
field.  The plaquette action below is fully nonlinear.  No assertion is made
that $R$ is a complete group-valued coarse connection or that it closes under
iteration.
